\documentclass[11pt,a4paper]{article}
\usepackage[margin=2cm,top=1.8cm,bottom=1.8cm]{geometry}
\usepackage[utf8]{inputenc}
\usepackage[T1]{fontenc}
\usepackage{hyperref}
\usepackage{xcolor}
\usepackage{amssymb}

\title{Inverting the Formalization Workflow:\\
  Prototyping an MPC Protocol in Rocq with an LLM}
\author{Cheng-Hui Weng}
\date{}

\begin{document}
\maketitle
\vspace{-1em}
\thispagestyle{empty}

\paragraph{The timeline problem.}
Formalization traditionally comes last.
A researcher studies the relevant mathematics, designs a construction informally, and only then translates the result into a proof assistant as a final verification step.
This ordering makes sense when the researcher already knows the mathematical landscape, but it imposes a high barrier when the construction requires unfamiliar domains.
What if formalization comes first---not as polished verification, but as a prototyping medium?
% REVIEW-FIX: [M5] footnote naming specific LLM on first body mention
An LLM\footnote{All LLM-assisted work reported here used Claude (Anthropic), specifically the Claude 3.5 Sonnet and Claude 3.5 Opus models.} can draft formal models from natural-language direction.
% REVIEW-FIX: [m1] "verdicts" as a verb is nonstandard; replace with "checks"
A dependent type checker checks each draft within minutes.
The human learns the domain from the structured output of each attempt: definitions that type-check reveal the right abstractions, and % REVIEW-FIX: [m2a] remove rhetorical "exactly"
failed proof obligations reveal which assumptions are missing.

\paragraph{The inverted workflow.}
The starting point was a hypothesis from separate ongoing work on generalizing secure multi-party computation (MPC): any surjective relation can serve as a secret-sharing mechanism if legitimate parties possess a trapdoor for reconstruction.
Covering spaces in algebraic topology match this pattern---the fiber over a basepoint is the secret, the monodromy action is the computation, and the covering map is the reconstruction.
Rather than studying group theory, algebraic geometry, and coding theory before attempting a protocol design, the author used formalization in Rocq/MathComp to test whether this connection holds.

From this seed, formalization becomes exploration.
The human does not need to master the required mathematical domains before starting.
The LLM contributes mathematical knowledge---definitions, lemma statements, proof strategies---that the human has never studied.
The type checker ensures each step is logically sound.
When a proof attempt fails, the failure is informative: % REVIEW-FIX: [m2b] remove rhetorical "exactly" (both occurrences)
a dependent-type mismatch reveals which mathematical structure is needed, and a universe inconsistency reveals which abstraction is wrong.
The human learns the domain from these failures and successes, iteration after iteration, with the type checker guaranteeing that each successful iteration stands on solid ground.

Three roles emerge.
The human provides judgment: which direction to explore, which abstractions to introduce, where to draw axiomatization boundaries, and when to abandon a dead end.
The LLM provides mathematical breadth: connecting group theory to combinatorics to algebraic geometry within a single proof development, contributing domain knowledge the human never formally studied.
% REVIEW-FIX: [M2] clarify that certainty is relative to stated axioms
% REVIEW-FIX: [M2] clarify certainty relative to axioms; [M4] inline contribution statement
The type checker provides certainty: if it compiles with zero \texttt{Admitted} lemmas, the logical structure holds relative to stated axioms.
This paper reports on this inverted workflow and the resulting SMC-PGG formalization (42 Rocq files, five mathematical domains, zero \texttt{Admitted}), with observations on the distinct roles of human judgment, LLM breadth, and type-checker certainty.

\paragraph{Case study: SMC-PGG.}
The seed idea grew into a four-layer formalization called SMC-PGG (Secure Multi-party Computation via Permutation-Group Graphs).

The \emph{protocol layer} formalizes a monodromy representation $\rho : G \to S_N$ mapping group elements to permutations on $N$ sheets, hiding a secret in the fiber structure of a covering space.
This is packaged as an HB mixin (\texttt{MonodromyReprType}) with session-typed programs for dealing, computing, and reconstructing.

The \emph{search-space layer} characterizes the adversary's view via Right-Angled Artin Group (RAAG) trace equivalence.
The Cartier-Foata theorem---proved axiom-free in Rocq---gives exact trace counts via a clique polynomial recurrence.
This was a domain the author had never encountered; its relevance was discovered through formalization when the proof required counting distinct computation histories that yield the same observable output.

The \emph{reconstruction layer} provides a threshold scheme interface instantiated by Massey's secret sharing from linear codes.
Hyperelliptic algebraic-geometry codes provide parametric $(k,T)$-thresholds via polynomial resultant computation, avoiding Riemann-Roch.
% REVIEW-FIX: [M1] "determines" overstates rigidity; ε is user-supplied, so use "constrains"
An algebraic rigidity framework (\texttt{AlgebraicRigidity} record) shows that one algebraic choice (a graph and its genus) constrains four protocol parameters, with $\varepsilon$ remaining user-supplied. A concrete star-graph instance and $\mathrm{PGL}(2,\mathbb{F}_q)$ bound are provided.

The \emph{security layer} establishes a collusion bound (% REVIEW-FIX: [C1] correct metric name to total variation distance
total variation distance $\leq \varepsilon + 2(T{-}1)/N$), Grover mitigation via ball-size enumeration, and an abelian collapse theorem showing that non-abelian groups are necessary for non-trivial security.

The formalization spans 42 Rocq files with zero \texttt{Admitted} lemmas, touching five mathematical domains: group theory, combinatorics, algebraic geometry, coding theory, and information theory.
The author entered the project knowing cryptographic protocol design.
The other four domains were learned through the formalization process, not before it.
Five properties of AG codes are axiomatized (Goppa weight bound, privacy surjectivity, monodromy compatibility, multiplicative closure, PGL automorphism bound), corresponding to classical algebraic geometry results that require Riemann-Roch machinery not yet available in any proof assistant.
This boundary is deliberate and principled---the protocol-algebraic layer is fully verified, and the axioms are standard textbook facts whose mathematical truth is not in question.

\paragraph{Observations.}
The workflow succeeded through a tight feedback loop: LLM drafts, type checker compiles or rejects within minutes, human steers.
Compilation success gave confidence in the mathematical structure; failure pointed precisely to missing invariants.
Limits remain: architectural decisions (\texttt{Record} vs \texttt{Section}, axiomatization boundaries) required human judgment, and dependent-type unification failures (e.g., rewriting under $\mathbb{Z}/(pq)\mathbb{Z}$) required human-level diagnosis.
The balance was productive: the human spent time on judgment and direction, not on tactic syntax or mathematical prerequisites.

% REVIEW-FIX: [M3] related-work paragraph with verified citations
\paragraph{Related work.}
LLM-assisted proving spans whole-proof generation \cite{first2023baldur}, informal-to-formal translation \cite{jiang2023dsp}, retrieval-augmented tactic search \cite{yang2023leandojo}, IDE integration \cite{song2024leancopilot}, and RL-based training with proof-checker reward \cite{alphaproof2024}.
On the cryptography side, EasyCrypt \cite{barthe2011easycrypt} provides machine-checked security proofs via probabilistic relational Hoare logic.
The present work differs in that the LLM is not optimized for proof search but serves as a domain-knowledge transfer mechanism for human-directed exploration.

\paragraph{Broader implication.}
This workflow makes mathematical applications more like software engineering, where developers learn libraries project by project, trusting type systems to catch misuse.
The LLM plays an analogous role: it surfaces relevant mathematical ingredients, while the type checker filters errors.
Formalization need not be the final step.
When an LLM serves as proof engineer and the type checker serves as quality gate, formalization becomes the first step, a prototyping medium that simultaneously validates ideas and teaches the domain.
The traditional order (learn $\to$ design $\to$ formalize) can be inverted (idea $\to$ formalize $\to$ learn) because the prover accepts no shortcuts.

% REVIEW-FIX: [M3] bibliography with all cited entries
\begingroup\footnotesize
\begin{thebibliography}{9}
\setlength{\itemsep}{0pt}\setlength{\parsep}{0pt}\setlength{\parskip}{0pt}

\bibitem{first2023baldur}
E.~First, M.~N.~Rabe, T.~Ringer, and Y.~Brun.
Baldur: Whole-proof generation and repair with large language models.
In \textit{ESEC/FSE 2023}, 2023.

\bibitem{jiang2023dsp}
A.~Jiang et al.
Draft, sketch, and prove: Guiding formal theorem provers with informal proofs.
In \textit{ICLR 2023}, 2023.

\bibitem{yang2023leandojo}
K.~Yang et al.
LeanDojo: Theorem proving with retrieval-augmented language models.
In \textit{NeurIPS 2023 (Datasets and Benchmarks)}, 2023.

\bibitem{song2024leancopilot}
P.~Song, K.~Yang, and A.~Anandkumar.
Lean Copilot: LLMs as copilots for theorem proving in Lean.
\textit{arXiv:2404.12534}, 2024.

\bibitem{alphaproof2024}
Google DeepMind.
AlphaProof and AlphaGeometry~2. Technical report, 2024.

\bibitem{barthe2011easycrypt}
G.~Barthe et al.
Computer-aided security proofs for the working cryptographer.
In \textit{CRYPTO 2011}, LNCS 6841, Springer, 2011.

\end{thebibliography}
\endgroup

\end{document}
